\section{Introduction}

\subsection*{The task}

The homework asks us to read the paper of Bryan and Leise~\cite{bryanleise}, to write our own code for the PageRank algorithm, to test it on the two webs of the paper (Figure~2.1 and Figure~2.2), to test it on the given dataset, and to solve two exercises of our choice.
We chose \textbf{Exercise~11} (a new page is added to the web and the ranking is computed with the matrix $\M$) and \textbf{Exercise~17} (how the damping factor $m$ should be chosen).

\subsection*{The dataset}

The dataset is the \emph{Hollins} web graph: 6012 pages of the website of Hollins University and 23\,875 links between them.
It is stored in the file \texttt{hollins.dat}.

\subsection*{Notation}

We use the notation of the paper.
A web has $n$ pages, numbered $1,\dots,n$.
The vector $\vx = (x_1,\dots,x_n)^T$ contains the importance scores, and we always scale it so that $\sum_i x_i = 1$.
$\Vone{\A}$ is the eigenspace of the matrix $\A$ for the eigenvalue~$1$.
When we quote a formula of the paper we use the paper's number, for example ``formula (2.1)'' or ``formula (3.2)''.
